\title{Parameter-Efficient Orthogonal Finetuning via Factorization}

\begin{abstract}
The prevalence of large foundation models continues to rise, yet the cost of training them from the ground up remains exceptionally high. Consequently, developing effective strategies for adapting these robust models to specific downstream tasks is of growing importance. In this work, we explore Orthogonal Finetuning (OFT)—a principled paradigm for downstream adaptation. While OFT is noted for its impressive generalization, it still requires the optimization of a substantial number of parameters due to the inherently high dimensionality of orthogonal matrices. To mitigate this, we approach OFT through the lens of information transmission and articulate several crucial criteria for enhancing parameter efficiency. Drawing inspiration from the Cooley-Tukey fast Fourier transform, known for its capacity-efficient information processing, we introduce an orthogonal parameterization scheme based on butterfly architectures. Integrating this into OFT results in a novel finetuning approach—Orthogonal Butterfly~(BOFT)—which offers heightened parameter efficiency. Notably, BOFT generalizes the OFT framework, which emerges as a particular instance of BOFT. We thoroughly evaluate the proposed method, BOFT, by adapting large vision transformers, extensive language models, and text-to-image diffusion models to diverse downstream applications spanning vision and language tasks.
\end{abstract}

\section{Introduction}

State-of-the-art models like ChatGPT~\cite{brown2020language,chen2021evaluating} and Stable Diffusion~\cite{rombach2022high} exemplify the outstanding generalization capabilities characteristic of modern large-scale foundation models. However, these advances are accompanied by a massive escalation in parameter counts—GPT-3, for instance, incorporates approximately 175 billion parameters. This explosion in scale renders training foundational models from scratch increasingly unattainable for most research groups. Therefore, meaningful progress in the field now hinges on methods for repurposing existing models rather than developing them anew. To this end, the efficient adaptation of foundation models to new tasks is imperative. Efficient model adaptation can be broadly categorized into three strategies: \emph{model finetuning}~\cite{hulora2022,zaken2022bitfit,guo2020parameter,raffel2020exploring,qiu2023controlling,zhang2022adaptive,chavan2023one}, which retains the original architecture but selectively updates portions of model parameters; \emph{adapter tuning}~\cite{rebuffi2017learning,houlsby2019parameter,pfeiffer2020adapterfusion,lin2020exploring,he2021towards}, which inserts new trainable components into the model; and \emph{prompt tuning}~\cite{li2021prefix,lester2021power}, which prepends learnable tokens to model inputs. Among these options, model finetuning is particularly attractive owing to its simplicity, efficacy, and the absence of any additional inference overhead.

The key idea behind model finetuning is to maintain a close resemblance between the pretrained and finetuned models according to specific metrics, thereby retaining the knowledge acquired during pretraining. A typical practice in contemporary finetuning techniques is to use a low learning rate, which heuristically limits the extent of changes between the original and updated models. Given the inherent structure of weight matrices, an alternative and more systematic strategy seeks to maintain the relational geometry within the weights, particularly the pairwise angular relationships among neurons. This rationale underpins the development of Orthogonal Finetuning~(OFT)~\cite{qiu2023controlling}, a new finetuning methodology in which an orthogonal matrix is uniformly applied to all neurons within a layer to rigorously preserve pairwise angles throughout finetuning. While OFT achieves favorable generalization and convergence outcomes for text-to-image diffusion model finetuning~\cite{qiu2023controlling}, the approach can suffer from a high parameter count due to the large size of the orthogonal matrices involved. To alleviate this, OFT incorporates a block-diagonal design, which decreases parameterization by restricting the transformation to independent submatrices. However, this gain in parameter efficiency introduces limitations: the enforced block-diagonal structure imposes a predetermined sparsity pattern, reducing the capacity of the orthogonal transformation to model diverse interactions—thereby restricting expressivity and hindering the approximation of canonical linear transformations. Moreover, determining an optimal sparsity configuration remains an unresolved issue.

The central challenge here lies in constructing a dense orthogonal matrix without incurring a prohibitive number of parameters. Although creating a dense $d$-dimensional orthogonal matrix typically necessitates $\mathcal{O}(d^2)$ parameters, we propose an alternative strategy: synthesizing the dense orthogonal matrix by composing several factorized sparse matrices. Our method stems from the realization that the parameter count can be diminished by exchanging greater computational complexity for reduced memory usage. This is made possible because the orthogonal matrix is expressed as a product of sparse factors, which requires repeated multiplication throughout the training process. To contextualize this matrix factorization, we recast the construction of a dense orthogonal matrix as an exercise in information propagation. In this paradigm, constructing the dense orthogonal matrix through a sequence of sparse multiplications is analogous to propagating information over a grid-like graph. This perspective opens up a range of design possibilities for sparse matrix factorizations that, while keeping the number of learned parameters small, still maintain the expressive capacity required to form dense matrices.

For improved parameter efficiency within our information transmission framework, we turn to butterfly networks, which are a foundational element in the Cooley-Tukey fast Fourier transform algorithm~\cite{cooley1965algorithm} and enable efficient data exchange between different nodes~\cite{klasing1994broadcasting}. The inherent butterfly topology utilized in the Cooley-Tukey approach suggests a specific method for sparse matrix decomposition, commonly known as butterfly factorization. Using this scheme, one can realize a dense matrix of dimension $d$ by sequentially multiplying $\mathcal{O}(\log d)$ sparse matrices, where each has just $\mathcal{O}(d)$ non-zero entries. By further constraining each sparse factor to be orthogonal, we establish a compact orthogonal parameterization involving only $\mathcal{O}(d\log d)$ parameters, marking a drastic decrease from standard full parameterizations. Building on this compact orthogonal structure, we introduce a novel parameter-efficient finetuning strategy named Orthogonal Butterfly~(BOFT). BOFT generalizes OFT as a particular case, furnishing a comprehensive framework for orthogonal finetuning. Both block-diagonal and butterfly matrix constructions exhibit a common feature: sparsity. This sparsity is harnessed to cut down the parameter footprint within orthogonal matrices. Notably, when compared with the low-rank mechanisms in LoRA~\cite{hulora2022}, our technique and LoRA both employ principal structured matrix families—namely, sparse and low-rank matrices~\cite{candes2011robust}—to promote parameter efficiency.

In contrast to the block-diagonal design used in OFT, which mediates between expressiveness and regularization, BOFT leverages the butterfly structure to achieve a more gradual and flexible interpolation between fully expressive matrices from the orthogonal group and highly regular identity matrices. This allows for the identification of a more compact set of candidates from the orthogonal group, tailored to the target applications. The butterfly structure, foundational in the implementation of efficient linear operators such as the discrete Fourier transform and the discrete cosine transform, motivates our hypothesis that this architecture naturally embeds an inductive bias that can enhance generalization and mitigate overfitting. The rationale behind this approach is rooted in the butterfly structure’s inherent capacity to reconstruct a variety of fundamental linear transforms—a property that the block-diagonal structure of OFT lacks. Our main contributions are as follows:

\begin{itemize}[leftmargin=*,nosep]
    \item We address parameter efficiency in orthogonal finetuning by introducing a novel information transmission perspective, reframing the challenge of generating a parameter-efficient dense orthogonal matrix as an information flow optimization over a graph with a grid topology.
    \item Building on butterfly patterns exemplified by the Cooley-Tukey algorithm, we introduce Orthogonal Butterfly, a parameter-efficient strategy for orthogonal finetuning developed within this information transmission paradigm.
    \item We offer theoretical justifications explaining how BOFT can drastically decrease the number of learnable parameters without compromising on model expressiveness or training robustness. Furthermore, due to its matrix factorization approach, BOFT inherently supports seamless weight interpolation (as demonstrated in Figure~\ref{fig:boft_interpolation}).
    \item For the first time, we extend orthogonal finetuning beyond the domain of controllable text-to-image generation~\cite{qiu2023controlling}, demonstrating its effectiveness as a versatile finetuning technique. Specifically, we evaluate BOFT across diverse adaptation tasks in both computer vision and natural language processing, where it not only surpasses existing state-of-the-art techniques in terms of parameter efficiency and generalization, but does so by a substantial margin.
\end{itemize}

\section{Related Work}

\textbf{Parameter-efficient finetuning (PEFT).} As the scale and capability of foundation models (\emph{e.g.}, \cite{brown2020language,radford2021learning,rombach2022high,kirillov2023segment}) continue to increase, the challenge of efficiently adapting these models to downstream applications has garnered significant attention~\cite{treviso2023efficient,ding2023parameter,lialin2023scaling}. The literature offers a wide array of PEFT techniques~\cite{houlsby2019parameter,aghajanyan2020intrinsic,hulora2022,edalati2022krona,wang2022adamix,gheini2021cross,zaken2022bitfit,guo2020parameter,sung2021training,ansell2022composable,lester2021power,li2021prefix,vu2022spot,he2021towards,mao2021unipelt,karimi2021compacter,liu2022few,sung2022lst,chen2022parameter,jia2022visual,chen2022adaptformer,zhang2022neural,jie2023fact,lian2022scaling,luo2023towards}, among which reparameterization-based strategies~\cite{aghajanyan2020intrinsic,hulora2022,edalati2022krona} are particularly relevant to the present study. In particular, LoRA~\cite{hulora2022} modifies pretrained weight matrices by appending a low-rank decomposition, consisting of two small matrices, which has yielded strong results for natural language processing tasks. While LoRA imposes a uniform rank constraint across all layers, subsequent work \cite{zhang2022adaptive,valipour2022dylora,zhang2023increlora} has introduced mechanisms for adapting the rank to each layer, thereby distributing the parameter resources more flexibly. This class of low-rank adaptation approaches has attained substantial popularity due to its conceptual simplicity and empirical effectiveness~\cite{dettmers2023qlora,zi2023delta,chavan2023one}. Motivated by insights into hyperspherical energy and its relation to model generalization~\cite{liu2018learning,liu2021orthogonal}, orthogonal finetuning (OFT) was recently introduced in \cite{qiu2023controlling} as an alternative that learns orthogonal transformations within a layer to reparameterize its neurons. OFT has demonstrated superior generalization capacity and more stable optimization dynamics than LoRA when adapting text-to-image diffusion models, although it typically involves more trainable parameters. Reducing the parameter footprint of OFT thus represents an important direction for research. Moreover, the potential applicability of OFT to a broader range of model adaptation tasks, beyond its initial focus on text-to-image diffusion~\cite{qiu2023controlling}, has not yet been explored. BOFT seeks to address this by leveraging butterfly factorization to enhance OFT's parameter efficiency, enabling the broader adoption of orthogonal finetuning for general adaptation purposes.

\textbf{Butterfly structures.} The radix-2 Cooley-Tukey algorithm decomposes the $N$-point discrete Fourier transform recursively into two subproblems of size $\frac{N}{2}$, naturally giving rise to a butterfly computational structure, which may be expressed as the product of several sparse matrices—collectively termed a butterfly matrix. Such matrices have been leveraged to represent orthogonal transformations for efficient computation, for instance, to eliminate the need for pivoting in Gaussian elimination~\cite{parker1995random}, to stabilize training dynamics in recurrent neural networks~\cite{jing2017tunable}, and for efficient kernel approximations~\cite{munkhoeva2018quadrature}. Recent advances~\cite{dao2019learning,dao2019kaleidoscope} have demonstrated that fast linear operators can be learned by parameterizing them with butterfly matrices. Butterfly representations also facilitate scalable and effective neural network training~\cite{chen2022pixelated,dao2022monarch}. In addition, applications of butterfly structures encompass fast matrix-vector multiplication~\cite{michielssen1996multilevel,o2010algorithm}, approximating large sparse matrices~\cite{li2015butterfly}, and improving data broadcast and network transmission protocols~\cite{klasing1994broadcasting,li2003linear,rajasingh2016transmission}. Unlike existing approaches, our work centers on deploying butterfly structures to augment the parameter efficiency of orthogonal finetuning methods.

\section{An Information Transmission View on Orthogonal Finetuning}

We begin by outlining key concepts in OFT. For fine-tuning a pretrained weight matrix, OFT reformulates the updated matrix as the multiplication of a learnable orthogonal matrix with the fixed, pretrained matrix. While LoRA accomplishes adaptation by incorporating an additive low-rank matrix to the pretrained weights, OFT instead employs a multiplicative, orthogonal matrix to realize updates. LoRA's approach to parameter efficiency is rooted in a low-rank formulation; conversely, original OFT achieves efficiency by structuring the orthogonal matrix in a block-diagonal fashion~\cite{qiu2023controlling}. The smaller the blocks along the diagonal, the greater the parameter reduction for OFT. Figure~\ref{fig:comp_lora} provides a visual comparison of these approaches. The use of orthogonal transformation for weight matrix fine-tuning is motivated by its ability to maintain pairwise angles among neurons~\cite{liu2018learning,liu2021orthogonal,qiu2023controlling}, thereby largely retaining the semantic knowledge acquired during pretraining. 

More specifically, OFT learns an orthogonal matrix $\bm{R}\in\mathbb{R}^{d\times d}$ for a given pretrained linear transformation $\bm{W}^0\in\mathbb{R}^{d\times n}$, changing the forward computation from $\bm{z}=(\bm{W}^0)^\top\bm{x}$ to $\bm{z}=(\bm{R}\bm{W}^0)^\top\bm{x}$, with input $\bm{x}\in\mathbb{R}^{d}$ and output $\bm{z}\in\mathbb{R}^{n}$. OFT accomplishes zero initialization by setting $\bm{R}$ as the identity matrix at the outset. To preserve the orthogonality of $\bm{R}$ during fine-tuning, Cayley parameterization~\cite{qiu2023controlling,liu2021orthogonal} is utilized: $\bm{R}=(\bm{I}+\bm{Q})(\bm{I}-\bm{Q})^{-1}$, where $\bm{Q}$ is skew-symmetric, i.e., $\bm{Q}=-\bm{Q}^\top$. Parameter efficiency is further attained by expressing $\bm{R}$ in a block-diagonal form: $\text{diag}(\bm{R}_1,\bm{R}_2,\cdots,\bm{R}_r)$, where each block $\bm{R}_i\in\mathcal{R}^{b\times b}$ for all $i$ and $br=d$. This block-diagonal design reduces parameters but imposes a limitation: it effectively partitions each neuron's dimensions—each corresponding to a column of $\bm{W}^0$—into $r$ groups, each being transformed independently by separate orthogonal matrices. Though empirically effective, assigning neuron dimensions to groups by index is somewhat arbitrary, highlighting the potential benefits of dense orthogonal matrices. This raises an important inquiry: \emph{Is it possible to design a dense orthogonal matrix while retaining parameter efficiency?}

To tackle this issue, we suggest constructing a dense orthogonal matrix as a product of several sparse orthogonal matrices. To establish a comprehensive and clear framework for analyzing the sparsity patterns in orthogonal matrix factorizations, we conceptualize the task of generating a dense orthogonal matrix in OFT as analogous to an information transmission process. In detail, forming a dense matrix $\bm{R}\in\mathbb{R}^{d\times d}$ through the multiplication of $m$ square matrices, given by $\thickmuskip=2mu \medmuskip=2mu \bm{R}=\bm{B}_m\bm{B}_{m-1}\cdots\bm{B}_1$, can be represented as the transmission of information across a $d\times (m+1)$ grid, as depicted in Figure~\ref{fig:info_trans}. This perspective is inspired by noting that a dense $d\times d$ matrix effectively embodies a full set of connections from one set of $d$ nodes to another. For matrix $\bm{R}$, a zero entry $R_{ij}$ signifies that information cannot flow from node $j$ to node $i$, whereas a nonzero entry means a possible information path. Thus, interpreting $\bm{R}$ as a product $\bm{B}_m\bm{B}_{m-1}\cdots\bm{B}_1$ equates to sequential exchanges of information, as determined by the graph structures generated by each $\bm{B}_i$. The flow of information proceeds in order, starting from $\bm{B}_1$ and ending with $\bm{B}_n$. As demonstrated in Figure~\ref{fig:info_trans}, we analyze the case of $\bm{R}=\bm{B}_5\bm{B}_4\bm{B}_3\bm{B}_2\bm{B}_1$, visualizing both the sparsity patterns and the corresponding induced graph. The graph is derived by expanding out the matrix product, treating each $\bm{B}_i$ as defining the connections from the $i$-th to the $(i+1)$-th layer of nodes. More concretely, entry $(j_1,j_2)$ in $\bm{B}_i$ encodes whether a directed connection exists from node $j_2$ at level $i$ to node $j_1$ at level $i+1$ (a zero indicates the absence of such an edge). For the composition $\bm{B}_5\bm{B}_4\bm{B}_3\bm{B}_2\bm{B}_1$ to yield a dense matrix, each node in the first level must be able to transmit to every node in the sixth level. Should we instead restrict attention to $\bm{R}=\bm{B}_4\bm{B}_3\bm{B}_2\bm{B}_1$—thus linking the first-level sources to fifth-level receivers—it becomes apparent that, for instance, information from node 1 cannot reach node 3. Consequently, the combined sparsity of $\bm{B}_4\bm{B}_3\bm{B}_2\bm{B}_1$ produces a zero at the $(3,1)$ entry. The same relationship between graph connectivity and matrix sparsity persists for $\bm{R}=\bm{B}_3\bm{B}_2\bm{B}_1$ and $\bm{R}=\bm{B}_2\bm{B}_1$. In summary, achieving a dense matrix $\bm{R}\in\mathbb{R}^{d\times d}$ requires that the $m$ constituent matrices $\bm{B}_m,\dots,\bm{B}_1$ together describe a set of directed edges across a $d\times (m+1)$ grid, where each edge links adjacent levels (columns), and ensures that each first-level node connects, possibly through multiple paths, to every node in the $(m+1)$-th level.

The information transmission perspective provides valuable insight into the sparsity configuration of the factorization matrices used in OFT. Figure~\ref{fig:blk_diag} illustrates the block-diagonal pattern exhibited by $\bm{R}$ in the standard OFT formulation. While this structure decreases the count of trainable parameters, it inherently limits our ability to achieve a fully dense matrix $\bm{R}$. Our objective is to synthesize a dense orthogonal matrix by composing $m$ sparse orthogonal matrices, minimizing the number of trainable parameters required in the process. Within the information transmission framework, two core principles guide our efforts: \emph{(i)} \textbf{dense connectivity}, ensuring that each node at the first level maintains at least one route to every node at the final level, and \emph{(ii)} \textbf{minimal free edges}, meaning that the total number of connections should be minimized subject to maintaining orthogonality. Importantly, the requirement of orthogonality introduces a nuanced restriction on the edges connecting consecutive levels. For instance, for each matrix $\bm{B}_i$ to possess full rank—an essential criterion for orthogonality—there must exist $d$ edges establishing a bijection between the nodes of the $i$-th and $(i+1)$-th levels. This condition implies that at least $d$ edges must connect adjacent levels (as seen in the illustration in Figure~\ref{fig:info_trans}, where this number is 4). These $d$ edges are essential for preserving orthogonality and are typically excluded from parameter counting, as their corresponding matrix elements are non-trainable (\emph{e.g.}, the non-zero entries in a $d\times d$ orthogonal matrix are fixed as $\pm 1$). Given that orthogonal matrices are parameterized by fewer variables compared to unrestricted matrices, the orthogonality constraint further intertwines the dependencies among edges. For example, in a $2\times 2$ orthogonal matrix, if the $(1,1)$ entry is zero, orthogonality requires that the $(2,2)$ entry is also zero (\emph{i.e.}, removing one edge enforces the removal of another). As such, the set of admissible connections is shaped by orthogonality, which may necessitate the inclusion or exclusion of certain edges. The nonzero structure of the resulting orthogonal matrix, when visualized, aligns well with the information transmission approach; in OFT, attention is primarily on the positions of nonzero, trainable elements in $\bm{R}$, with their magnitudes being of lesser consequence.

A straightforward dense linkage between two layers requires $\mathcal{O}(d^2)$ connections (i.e., a full dense orthogonal matrix), which results in $d^2-d$ individual connections (for instance, with $d=4$, there are 12 edges). As demonstrated in Figure~\ref{fig:info_trans}, a particular example of feasible matrix factorization uses only $10$ edges in total, a reduction compared to a dense orthogonal matrix. This approach offers a graphical lens to analyze the parameter efficiency of OFT and supports the construction of feasible factorizations within this framework. Drawing on the structure of the butterfly graph, originally from the Cooley-Tukey algorithm~\cite{cooley1965algorithm}, we note that such topologies are capable of densely interconnecting $d$ input and $d$ output nodes with only $\mathcal{O}(d\log d)$ edges. For comparison, the topology depicted in Figure~\ref{fig:info_trans} requires $10$ edges for full connectivity, whereas a butterfly network can achieve the same with just $8$ edges. We will now describe how leveraging the butterfly structure can further enhance parameter efficiency.

\section{Orthogonal Parameterization by Butterfly Factorization}

The butterfly configuration, central to the Cooley-Tukey algorithm, facilitates rapid Fourier transform computation. The key property in the Fourier transform is that a localized modification in the frequency domain propagates to affect the entire spatial domain—mirroring the core requirement in our information transmission context: enabling every first-level node to communicate with all nodes at the final level. Furthermore, the butterfly network has also gained popularity as a computer network architecture~\cite{solihin2015fundamentals,li2011linear} for high-efficiency data transfer. Let $k\geq 2$ denote a power of $2$; we define the \emph{butterfly factor} $\bm{B}^F(k)$ by

\begin{equation}\label{eq:but_factor}
\begin{aligned}
    \bm{B}^F(k)=\begin{bmatrix}
    \text{diag}(\bm{d}_1) & \text{diag}(\bm{d}_2) \\
    \text{diag}(\bm{d}_3) & \text{diag}(\bm{d}_4)
    \end{bmatrix}\in\mathbb{R}^{k\times k},
\end{aligned}
\end{equation}

with each $\bm{d}_i\in\mathbb{R}^{\frac{k}{2}}$ for $i=1,2,3,4$. Given $d=2^N$, we subsequently define the $d$-dimensional \emph{butterfly matrix} $\bm{B}(d)\in\mathbb{R}^{d\times d}$ recursively as follows:

\begin{equation}
\begin{aligned}
    \!\!\bm{B}(d)=\tilde{\bm{B}}(d,d)\!\cdot\!\begin{bmatrix}
    \bm{B}_1(\frac{d}{2})\!\!\!\!\! & \bm{0} \\
    \bm{0}\!\!\!\!\! &\bm{B}_2(\frac{d}{2})
    \end{bmatrix}= \tilde{\bm{B}}(d,d)\tilde{\bm{B}}(d,\frac{d}{2})\cdots\tilde{\bm{B}}(d,2),
\end{aligned}
\end{equation}

where $\bm{B}_1(\frac{d}{2})$ and $\bm{B}_2(\frac{d}{2})$ represent two butterfly matrices, each of dimension $\frac{d}{2}$. We introduce the notion of a \emph{butterfly component}, given by $\thickmuskip=2mu \medmuskip=2mu \tilde{\bm{B}}(d,k)=\text{diag}(\bm{B}^F_1(k),\cdots,\bm{B}^F_{d/k}(k))$, which forms a block-diagonal matrix of size $d\times d$, with each block of size $k$. Here, each $\bm{B}^F_i(k)$, for all $i$, is a butterfly factor as specified in Equation~\ref{eq:but_factor}. With these definitions, we proceed to use butterfly matrices to parameterize an orthogonal matrix. To realize this, it suffices to ensure that all block-diagonal factors $\tilde{\bm{B}}(d,k)$ comprising the butterfly matrix $\bm{B}(d)$ are orthogonal. To illustrate, consider the block-diagonal matrix $\tilde{\bm{B}}(d,2)$ with block size 2; each block can be parameterized to be orthogonal using the Cayley transform or a $2$-dimensional rotation, following the approach of \cite{liu2021orthogonal,qiu2023controlling}. The sparsity structure of any butterfly component corresponds to a permutation of the non-zero structure in $\tilde{\bm{B}}(d,2)$; consequently, every butterfly component can also be parameterized to be orthogonal with ease. This approach constructs an efficient parameterization of orthogonal matrices by assembling multiple $2\times 2$ orthogonal matrices. Following \cite{chen2022pixelated}, we further generalize butterfly matrices and define a block butterfly component $\tilde{\bm{B}}^b(d,k)$ where each element $\bm{d}_i$ (for all $i$) is promoted to a $b\times b$ matrix. To preserve orthogonality for the block butterfly component $\tilde{\bm{B}}^b(d,2)$, each $2b\times 2b$ block is likewise parameterized as an orthogonal matrix. The nonzero structure in higher-order butterfly components $\tilde{\bm{B}}^b(d,k)$ for $k>2$ amounts to block-wise permutations of the nonzero pattern in $\tilde{\bm{B}}^b(d,2)$; as such, orthogonal parameterization applies analogously. All together, the BOFT forward pass is given by

\begin{equation}
    \begin{aligned}\nonumber
    \bm{z}=\big(\bm{R}(m,b)\cdot\bm{W}^0\big)^\top\bm{x},~~~~\text{s.t.}~\bigg{\{}\bm{R}(m,b)=\prod_{i=1}^m \tilde{\bm{B}}^b_{(d,i)}~~\&~~\underbrace{\big(\tilde{\bm{B}}^b_{(d,j)}\big)^\top\tilde{\bm{B}}^b_{(d,j)}=\tilde{\bm{B}}^b_{(d,j)}\big(\tilde{\bm{B}}^b_{(d,j)}\big)^\top\!=\bm{I}_d}_{\forall j\in [1, m]}\bigg{\}},
    \end{aligned}
\end{equation}

In this context, we use the notation $\tilde{\bm{B}}^b(d,2^{m - i+1})$ as a shorthand for $\tilde{\bm{B}}^b_{(d,i)}$, and $\bm{I}_d$ indicates the $d\times d$ identity matrix. The orthogonal transformation $\bm{R}(m,b)\in\mathbb{R}^{d\times d}$ is constructed as a sequence of multiple butterfly-based orthogonal matrices. For clarity, we refer to BOFT parametrized with $\bm{R}(m,\frac{b}{2})$ as BOFT($m$,$b$), where $b$ must be at least 2. When setting $m=1$, BOFT($1$,$b$) simplifies to a block-diagonal variant of OFT~\cite{qiu2023controlling}, with blocks of size $b$. If $b$ equals $d$, then BOFT($1$,$d$) matches the full dense version of OFT~\cite{qiu2023controlling} with no block restrictions. By choosing $m=\log \frac{2d}{b}$, BOFT($m$,$b$) configures $\bm{R}$ to be a block butterfly matrix $\bm{B}^{\frac{b}{2}}(d)$, yielding a dense orthogonal $\bm{R}$. More broadly, BOFT($m$,$b$) utilizes $\frac{1}{2}(b-1)dm$ learnable parameters to adapt a $d\times n$ linear mapping. Employing butterfly matrices specifically ($m=\log d, b=2$), BOFT leverages $\mathcal{O}(d\log d)$ parameters. In contrast, a conventional OFT with a complete orthogonal matrix has parameter complexity $\mathcal{O}(d^2)$, while the block-diagonal OFT with $r$ blocks demands $\mathcal{O}(bd)$. Consequently, to realize a dense orthogonal transformation, standard OFT requires $b=d$ for the block size, whereas BOFT is flexible in its choice of $b$.

\textbf{BOFT Identity Initialization}. Typical parameter-efficient fine-tuning strategies begin with the weights of the original pre-trained network, minimizing the divergence of the new model from its starting point. For instance, LoRA adopts zero initialization for its rank-reduced parameter matrices. In the BOFT framework, all butterfly components are initialized as identity matrices (specifically, their underlying skew-symmetric matrices in the Cayley representation are zero-initialized).

\textbf{Multiplicative Dropout for BOFT}. To reduce overfitting, LoRA~\cite{hulora2022} introduces a Dropout operation on the low-rank weight deltas. Standard Dropout~\cite{srivastava2014dropout} integrates seamlessly into LoRA’s additive update scheme, but it is incompatible with the multiplicative structure of BOFT’s parameterization. To address this challenge, we introduce a multiplicative Dropout adapted for BOFT. Recognizing that $\bm{R}(m,b)$ consists of $m$ orthogonal butterfly sub-matrices, each of which can be reorganized as $2b\times 2b$ block-diagonal orthogonal matrices, our method randomly selects $p_1$ percent of butterfly sub-matrices and, within each, selects $p_2$ percent of their diagonal blocks. These chosen blocks are then substituted with identity matrices during training.

\section{Intriguing Insights and Discussions}

\textbf{Expressivity of BOFT}. The butterfly structure, especially when combined with permutations, is capable of exactly reconstructing many well-known fast linear transforms~\cite{dao2019learning,dao2019kaleidoscope} (\emph{e.g.}, fast Fourier transform, Hadamard transform). However, the capability of our orthogonal butterfly matrix to approximate an arbitrary orthogonal matrix has not been fully characterized. To investigate this, we carry out experiments to approximate a randomly generated dense orthogonal matrix~\cite{anderson1987generation} of dimension $512\times 512$. The outcomes, depicted in Figure~\ref{fig:para_eff}, are averaged over 10 independent random initializations. The y-axis displays the approximation error, while the x-axis corresponds to the count of effective trainable parameters. Each colored curve represents BOFT with a given block size, and the leftmost marker reflects the performance of BOFT($1$,$b$), i.e., the standard block-diagonal OFT with block size $b$. Generally, BOFT demonstrates higher parameter efficiency compared to OFT. For instance, BOFT($9$,$2$) achieves greater expressivity than BOFT($1$,$16$) despite employing significantly fewer parameters. Configurations with smaller block size $b$ and larger $m$ tend to deliver higher parameter efficiency; as an example, BOFT($6$,$4$) attains comparable approximation errors to BOFT($2$,$16$) but requires far fewer parameters. In summary, the butterfly matrix encompasses a more structured subclass of the orthogonal group than the block-diagonal configuration, resulting in BOFT having a provably superior expressive power over OFT for equivalent block sizes.

\begin{theorem}[Expressivity of BOFT]\label{thm:exp_boft}
    BOFT surpasses OFT in expressiveness when utilizing the same block size. In order for a butterfly matrix to represent all orthogonal matrices of dimension $d$, it suffices to form products of the form $\bm{B}_{d-1,1}(d)\bm{B}^\top_{d-1,2}(d)\cdots\bm{B}_{1,1}(d)\bm{B}^\top_{1,2}(d)$, where each $\bm{B}_{i,j}(d)$ is a butterfly matrix for all relevant $i$ and $j$.
\end{theorem}

Theorem~\ref{thm:exp_boft} points towards a straightforward extension of BOFT, in which the ultimate orthogonal matrix is defined as $\thickmuskip=2mu \medmuskip=2mu \bm{R}^G(m_1,b_1,m_2,b_2,l)=\bm{R}_{l,1}(m_1,b_1)\bm{R}_{l,2}^T(m_2,b_2)\cdots\bm{R}_{1,1}(m_1,b_1)\bm{R}_{1,2}^T(m_2,b_2)$, where $\bm{R}_{i,j}^T(m,b)$ is an orthogonal matrix used within the BOFT framework. Specifically, setting $m_1=m_2=\log d$, $b_1=b_2=2$, and $l=d-1$, allows $\bm{R}^G(m_1,b_1,m_2,b_2,l)$ to span the entire orthogonal group. This matrix composition is referred to as the kaleidoscope hierarchy~\cite{dao2019kaleidoscope}. Nevertheless, greater expressivity does not inherently ensure superior fine-tuning outcomes; for example, while full finetuning is universally expressive, it frequently underperforms in practice. The crucial factor for effective model finetuning lies in balancing expressivity with regularization. BOFT increases the parameter space during finetuning by introducing structural priors, thereby facilitating a more optimal trade-off between model expressivity and regularity.

\textbf{Spectral properties}. Compared to LoRA, orthogonal finetuning typically results in superior spectral characteristics, as it exactly maintains the spectral norm of the original pretrained matrix $\bm{W}^0$. This is evident through its singular value decomposition: expressing $\bm{W}^0$ as $\bm{U}\bm{\Sigma}\bm{V}^\top$, where $\bm{U}$ and $\bm{V}$ are orthogonal matrices and $\bm{\Sigma}$ contains the singular values along its diagonal. Both OFT and BOFT operate by left-multiplying with an orthogonal matrix $\bm{R}$, producing updated weights of the form $\bm{R}\bm{U}\bm{\Sigma}\bm{V}^\top$, which leaves the maximal singular value (i.e., the spectral norm of $\bm{W}^0$) unchanged. Maintaining the spectral norm in this manner has been empirically linked to enhanced training robustness and improved generalization~\cite{miyato2018spectral,yoshida2017spectral}. Additional mathematical insights are detailed in Appendix~\ref{app:sn_property}.

\textbf{Orthogonal finetuning as learning bilinear similarity}. The formulation of BOFT can be interpreted as optimizing the bilinear similarity function $\bm{w}^0_i\bm{R}\bm{x}$, where $\bm{w}^0_i$ refers to the $i$th column of the weight matrix $\bm{W}^0$. Here, BOFT is conceptualized as learning a bilinear similarity matrix $\bm{R}$ subject to a strict orthogonality constraint, thereby establishing a clear relationship to both distance metric learning~\cite{xing2002distance} and bilinear forms~\cite{roman2005advanced}.

\textbf{Inductive bias and generalization in BOFT}. In BOFT, $\bm{R}(m,b)$ is generally chosen from a highly structured subset of the orthogonal group, thereby restricting the hypothesis space and inducing an inductive bias. The structure introduced by butterfly factorization is argued to be advantageous for generalization, owing to its shared structure with several fundamental linear transforms such as the discrete Fourier transform, discrete sine/cosine transform, and the Hadamard transform~\cite{dao2019kaleidoscope}. Furthermore, the inherent sparsity of the matrix factorization utilized in BOFT may also contribute additional forms of implicit inductive bias~\cite{gunasekar2017implicit,li2018algorithmic,liu2019neural}.

\textbf{Comparison to butterfly-based sparse training}. Several studies~\cite{dao2019learning,dao2019kaleidoscope,chen2022pixelated,dao2022monarch} have explored sparse training through butterfly parameterizations. These approaches generally involve directly reparameterizing weight matrices with butterfly structures and training neural networks from the beginning. In contrast, \cite{dao2022monarch} investigates finetuning by first projecting pretrained weights onto a modified butterfly matrix before optimizing these components for downstream applications. BOFT introduces a markedly distinct finetuning method, in which transformation is achieved via layer-wise shared weight matrices.

\input{rebuttal-results}

\section{Concluding Remarks and Limitations}

This work introduces Orthogonal Butterfly, a broadly applicable and parameter-efficient finetuning technique for foundation models rooted in the butterfly structure. Our central observation is that improved parameter efficiency is attainable by expressing a dense orthogonal matrix as a product of multiple sparse orthogonal matrices. To streamline the search for suitable factorizations, we establish a graphical information transmission perspective. Through this lens, we identify that the butterfly pattern serves as an effective strategy for constructing sparse orthogonal matrix factorizations. We provide empirical results demonstrating the efficacy of BOFT in finetuning large-scale models in language, vision, and text-to-image generation domains, substantiating BOFT’s advantages as a versatile finetuning method.

Nonetheless, BOFT is not without shortcomings. Because the final orthogonal matrix in BOFT is a composition of several orthogonal matrices, training incurs a moderate increase in runtime compared to OFT, and addressing this computational cost is still an unresolved question. However, after finetuning, the orthogonal matrices obtained by BOFT can be directly merged into the pretrained model, ensuring there is no extra latency during inference. Furthermore, it remains unclear whether butterfly networks are the optimal design for information propagation. Our information transmission framework suggests possible inspiration from the field of computer networking, where questions regarding the efficiency of network topologies are well-studied. We anticipate that even more effective network structures could be harnessed to build dense orthogonal matrices.

\end{document}